\PassOptionsToPackage{numbers,sort&compress}{natbib}
\documentclass{article}

% paper_public.tex defines \PublicVersion before loading this file.
\providecommand{\PublicVersion}{0}
\ifnum\PublicVersion=1\relax
  \usepackage[preprint]{neurips_2026_vericode}
\else
  \usepackage{neurips_2026_vericode}
\fi
\workshoptitle{AI for Verifiable Coding}

\usepackage{amsmath,amssymb}
\usepackage{microtype}
\usepackage{booktabs,tabularx,array}
\usepackage{float}
\usepackage{graphicx}
\usepackage{xcolor}
\usepackage{listings}
\usepackage{listingsutf8}
\usepackage{tikz}
\usepackage[colorlinks=true,allcolors=blue]{hyperref}
\usepackage[nameinlink]{cleveref}
\usepackage[T1]{fontenc}
\usepackage[utf8]{inputenc}

% The review style loads lineno only for the anonymous submission build.
% Keep the historical theorem display usable in the public preprint as well.
\ifnum\PublicVersion=1\relax
  \newenvironment{nolinenumbers}{}{}
\fi

\newcommand{\Lean}{\textsc{Lean}}
\newcommand{\DK}{Davis--Kahan}
\input{generated/practitioner_account_macros.tex}
\input{generated/evidence_macros.tex}

% Set \PublicVersion to 1 in paper_public.tex for author names and public URLs.
\newif\ifshowartifacturl
\ifnum\PublicVersion=1\relax
  \showartifacturltrue
\else
  \showartifacturlfalse
\fi
\newcommand{\artifacturl}{https://github.com/AIQ-Kitware/aiq-dkps-formalization}

\definecolor{leankeyword}{HTML}{243B6B}
\definecolor{leantactic}{HTML}{70477A}
\definecolor{leantype}{HTML}{1F6473}
\definecolor{leancomment}{HTML}{39734D}
\definecolor{leanstring}{HTML}{8A4B2A}
\definecolor{leannumber}{HTML}{777777}
\definecolor{leanbackground}{HTML}{F7F7F3}
\lstdefinelanguage{lean}{%
  sensitive=true,
  morecomment=[l]{--},
  morecomment=[s]{/-}{-/},
  morestring=[b]",
  morekeywords=[1]{import,open,namespace,section,end,noncomputable,variable,variables,
    universe,universes,set_option,attribute,axiom,opaque,theorem,lemma,def,abbrev,
    example,instance,structure,class,inductive,notation,local,deriving,extends,
    by,where,let,have,show,fun,match,with,if,then,else,do,return,in,forall},
  morekeywords=[2]{simp,simpa,rw,erw,dsimp,unfold,ring,ring_nf,nlinarith,linarith,
    positivity,norm_num,field_simp,omega,exact,refine,apply,assumption,intro,intros,
    ext,funext,cases,rcases,obtain,constructor,left,right,classical,by_contra,
    contrapose,convert,using,suffices,calc},
  morekeywords=[3]{Prop,Type,Nat,Int,Real,Fin,Set,Submodule,Matrix,EuclideanSpace,
    LinearMap,ContinuousLinearMap,SymmetricNormingFunction,IsSelfAdjoint,Reduces},
}
\lstdefinestyle{leanpaper}{%
  language=lean,
  columns=fullflexible,
  backgroundcolor=\color{leanbackground},
  basicstyle=\ttfamily\footnotesize,
  commentstyle=\color{leancomment}\itshape,
  keywordstyle=[1]\color{leankeyword}\bfseries,
  keywordstyle=[2]\color{leantactic},
  keywordstyle=[3]\color{leantype},
  stringstyle=\color{leanstring},
  numberstyle=\tiny\color{leannumber},
  breakatwhitespace=false,
  breaklines=true,
  captionpos=b,
  keepspaces=true,
  numbers=none,
  showspaces=false,
  showstringspaces=false,
  showtabs=false,
  tabsize=2,
  aboveskip=0.6em,
  belowskip=0.6em,
  % This deliberately follows the MCC paper's \Real/\root pattern: the
  % listings input contains ASCII stand-ins and literate= does the typesetting.
  % IMPORTANT: the left-hand sides below are ASCII presentation sentinels,
  % not literal Lean Unicode. scripts/build_evidence.py writes the exact
  % extracted Lean signature to generated/historical_scope_mismatch_exact.lean,
  % then writes a separate listings input in which only display-sensitive Lean
  % notation is replaced by these sentinels. listings resolves each sentinel to
  % LaTeX here. Keep this indirection: it avoids Unicode decoding/editability
  % problems in listings/Overleaf while preserving an exact Lean sidecar for
  % audit. Do not replace the sentinel keys below with literal Lean glyphs.
  literate=
    {\\LeanLitCLMapC}{{$\mathrel{\to}\!\mathtt{L}[\mathbb{C}]$}}5
    {\\LeanLitComplex}{{$\mathbb{C}$}}1
    {\\LeanLitReal}{{$\mathbb{R}$}}1
    {\\LeanLitLe}{{$\le$}}1
    {\\LeanLitSubsetEq}{{$\subseteq$}}1
    {\\LeanLitForall}{{$\forall$}}1
    {\\LeanLitMem}{{$\in$}}1
    {\\LeanLitOrth}{{$^{\perp}$}}1
    {\\LeanLitDisj}{{$\lor$}}1
    {\\LeanLitConj}{{$\land$}}1
}
\crefname{lstlisting}{formalization}{formalizations}
\Crefname{lstlisting}{Formalization}{Formalizations}
\renewcommand{\lstlistingname}{Formalization}

\title{Did We Really Formalize Davis--Kahan? Semantic Alignment for LLM-Assisted Lean Formalization}
\author{%
  Jonathan Crall \quad Brian Hu \\
  Kitware, Inc.
  \And
  Edward Wang \quad Carey E. Priebe \\
  Johns Hopkins University
}
\date{}

\begin{document}
\maketitle

\begin{abstract}
We report an LLM-assisted Lean formalization of the 29 mathematical results
tracked from Davis and Kahan's 1970 paper.  The project is a case study in
semantic alignment: whether machine-checked formal statements preserve the
mathematical claims made in the source.  Lean verifies proofs of encoded
propositions, while source correspondence still requires checking hypotheses,
objects, and scope.  The current development contains checked treatments for
all 29 tracked results, including a counterexample and repaired theorem for
Proposition~4.4.  We describe the formalization and source-review process,
examples of discrepancies found during development, and
\PractitionerAccountCount{} public first-person accounts of AI-assisted Lean
work.  The case study illustrates why successful compilation alone does not
establish that an existing theorem has been formalized as intended.
\end{abstract}

\section{Introduction}

Lean checks that a proof term establishes the proposition encoded in a theorem
statement \citep{demoura2021lean4,mathlib2020}.  Establishing that the proposition
expresses a source theorem also requires checking its hypotheses, objects, and
scope.  This agreement is called \emph{semantic alignment}
\citep{lu2025formalalign}.  For a research paper, the intended claim includes
definitions, inherited assumptions, and surrounding mathematical context.
Chow illustrates the risk with an agent that replaces ``all doohickeys are
excellent'' by ``all special doohickeys are quasi-excellent''
\citep{chow2026mindreading}.  A valid proof of the replacement would leave the
intended theorem unresolved.

We encountered this problem while formalizing Davis and Kahan's 1970
spectral-perturbation results \citep{davis1970rotation}.  We wanted reusable
bounds for downstream applications, including the paper's real and complex
Hilbert-space settings.  Source review exposed both formal statements that
covered less than the source and a false printed proposition.  This paper
examines how we organized that review when LLMs could produce more Lean than
we could inspect line by line.

\Cref{fig:workflow} summarizes the process that developed around these
source-correspondence failures.  The source material remains visible through
context gathering, decomposition, foundation search, formalization, and a
separate source-to-statement review.  Section~\ref{sec:system} describes these
stages in more detail.

\begin{figure}[H]
  \centering
  \includegraphics[width=0.98\linewidth]{figures/formalization_workflow.png}
  \caption{Formalization workflow used in this case study.  Failed semantic
  review returns a target to the loop; accepted results contribute checked
  theorems and reusable foundations.}
  \label{fig:workflow}
\end{figure}

\section{Related Work}
\label{sec:related}

Using LLM-assisted web search, we collected \PractitionerAccountCount{} public
first-person project/workflow accounts of AI-assisted Lean work.  Multiple public
sources can support one account, and accounts known to share a practitioner are
marked explicitly; the \PractitionerAccountCount{} rows form
\PractitionerOverlapClusterCount{} practitioner-overlap clusters, not a count of
unique people.  They describe different ways of deciding whether generated Lean
expresses the intended mathematics.  The collection is not necessarily
representative; Appendix~\ref{app:data} describes the source records,
qualifications, and screening log.  Collins et al. study initial human--AI
formalization workflows, while Ilin and Nugent compare a sorry-free
semi-autonomous development before and after expert review
\citep{collins2026workflows,ilin2026sorries}.

Kahle and Saiki describe checking definitions, hypotheses, and decomposition
while agents write much of the proof code
\citep{kahle2026keeptrying,saiki2026double}.  Deng and Shum separate successful
builds from source-fidelity review; they report an initial variance declaration
that compiled while imposing weaker domain discipline than the textbook
intended \citep{deng2026probability}.  De Dios Pont et al. report concentrating
intensive review on theorem statements and definitions, then returning to
understand and generalize an argument after obtaining a Lean certificate
\citep{dedios2026banach}.  Other accounts put less emphasis on reading generated
Lean directly: Suomela studies the mathematical proof first, Gowers reports
using Aristotle without needing Lean knowledge, and Enomoto describes a later
``proof digestion'' phase
\citep{suomela2026workflow,gowers2026leiden,enomoto2026essay}.

Other accounts use a second model to mediate review.  Nowak and F\"ugger used
Claude to check Codex formalizations for added hypotheses, weakened statements,
and proof holes \citep{nowak2026workflow}.  Kova\v{c} asked another model to
translate Lean into English because he could not read it directly
\citep{kovac2025erdos189}; Ilin reports adversarial review among multiple systems
\citep{ilin2026vml}.  Douglas et al. cross-validated conjectured helper results
and proof plans across several models \citep{douglas2026qft}.  Human
counterexamples also catch statement errors: Zhang et al. report three
misformalized targets that survived repeated AI self-judgment before human
counterexamples exposed them \citep{zhang2026slt}; Saiki and Oum report related
uses of counterexamples and source correction
\citep{saiki2026double,oum2026preprint}.

The expanded corpus also includes software-facing trust boundaries.  Yadav
reports carefully checking every definition and theorem statement while largely
ignoring generated proof implementations; Warren reports five AI-written formal
statements that type-checked and passed numerical validation before proof
obligations exposed defects; and Zaliva reports parser and semantics defects, a
human-directed semantic redesign, and an initially wrong progress property
\citep{yadav2026verification,warren2026tessera,zaliva2026semantics}.  A verified
polygon-intersection project instead makes a short human-reviewed specification
the trust boundary while leaving AI-written implementation and proof files
unread \citep{schildep2026polygon}.

The added mathematical accounts show a second range of supervision.  Armstrong
and Kempe report supplying detailed proof blueprints and architecture while
agents wrote the Lean; Aslyan describes selecting symbols and redirecting
modeling decisions while the agent implemented them.  Vergo, Bernier, and
Schmitt explicitly solicit review of whether generated formalizations encode
the intended mathematics, while Davis preserves the blueprint, raw model output,
cleanup diff, and pinned environment for independent inspection
\citep{armstrong2026degiorgi,aslyan2026friedberg,vergo2025aiproofs,
bernier2026classfield,schmitt2025extremal,davis2026gdlean}.

Within the \PractitionerAccountCount{} accounts, \AccountSemanticMismatchCount{}
explicitly report a statement, definition, scope, or correspondence mismatch,
with \AccountSemanticMismatchQualifiedCount{} additional qualified cases.
\AccountSeparateAIReviewCount{} describe a separate AI review role, and
\AccountPersistentStateCount{} keep substantive project state outside the chat.
Formalization exposed a source defect in \AccountSourceDefectCount{} accounts;
\AccountCounterexampleCount{} explicitly used a counterexample to diagnose or
communicate a problem.  These categories can overlap and describe only the
accounts collected here.

In a controlled 100-agent experiment, Paglieri et al. report agents using local
notation to change how unchanged benchmark text elaborated while still passing
a submission harness based on template checks and Lean compilation
\citep{paglieri2026swarm}.  Lean checked the submitted propositions; the harness
did not ensure that the benchmark text retained its intended elaboration.

Semantic-review tools address several parts of this process.  FormalAlign
evaluates informal/formal agreement \citep{lu2025formalalign}; ShadowBench tests
formal statements through auxiliary statements and implication checks
\citep{han2026shadowbench}.  EconCSLib presents source text, Lean statements,
mathematical renderings, and human match/mismatch judgments together
\citep{garg2026econcs}.  Lean Atlas and Lean Compass prune proof-only dependencies
to expose definitions that can affect a theorem's meaning
\citep{yanahama2026leanatlas}.  LeanMarathon uses adversarial target review,
FormaTheoria supports independent review and source reconciliation, and
LeanArchitect synchronizes informal blueprint metadata with Lean declarations
\citep{zhang2026leanmarathon,nie2026formatheoria,zhu2026leanarchitect}.
These systems provide context for our result register and review interface.

Compiler access is supported by LeanDojo, Pantograph, and LeanInteract
\citep{yang2023leandojo,aniva2024pantograph,poiroux2025leaninteract}.
As this project grew, we developed smaller project-specific tools for querying
compiler-derived declarations and maintaining the source-to-statement records;
the supplementary artifact contains those tools and their exact interfaces.

\section{What we formalized}
\label{sec:dk}

Davis and Kahan bound the change in invariant subspaces under self-adjoint
perturbations.  In matrix terms, let $A=A^*$ (i.e., $A$ is self-adjoint) and let
$E_0$ have orthonormal columns spanning a \emph{candidate subspace} $V$, which
is to be compared with an exact invariant subspace.  The candidate-space
\emph{compression}, $A_0=E_0^*AE_0$, is obtained by using $E_0$ to embed
candidate coordinates into the ambient space, applying $A$, and using $E_0^*$
to project the result back to candidate coordinates.  Writing $P_V=E_0E_0^*$ for the orthogonal projector
onto $V$, the \emph{residual}
\[
  R=AE_0-E_0A_0=(I-P_V)AE_0
\]
is exactly the part of $AE_0$ that lies outside $V$; it vanishes exactly when
$A$ preserves the candidate subspace.  Principal angles measure the angles
between corresponding subspace directions.  With spectral separation $\delta>0$,
one directed
$\sin 2\Theta$ result has the form
\begin{equation}
  \delta\,N(\sin 2\Theta_0) \le 2N(R),
  \label{eq:sin2theta}
\end{equation}
where $N$ is a unitarily invariant norm and $\Theta_0$ denotes the directed
angle from the candidate to the exact subspace.  Here the gap separates the
spectrum on the exact subspace from that on its orthogonal complement.
The source also covers real and complex separable Hilbert spaces, including
unbounded self-adjoint operators under its domain and bounded-residual
conditions \citep{davis1970rotation}.  Appendix~\ref{app:formalization1} gives a
reader's guide to the mathematics and the Lean signature below.

Our 29 targets comprise the four unnumbered headline theorems in Section~2 and
every named theorem, proposition, lemma, and corollary established in the paper.
Appendix~\ref{app:data} specifies the counting convention.  At the manuscript
snapshot, all 29 have checked treatments.  Source-to-signature review has
continued to refine some source-facing theorem boundaries, including the norm
domain of the sine-theta theorem discussed below.  We therefore report the
current review state without treating it as a certificate of complete semantic
correspondence.

The development required projection and principal-angle geometry,
singular-subspace geometry, unitarily invariant norms and operator ideals,
Borel spectral calculus, unbounded self-adjoint operator theory, real/complex
transfer, direct rotations, and Sylvester equations.  Parts of the spectral and
operator-theoretic infrastructure were copied, adapted, or generalized from the
Spectra formalization \citep{bornemann2026spectra}.  Reusable components developed
here are being prepared for Tau Ceti, an AI-authored Lean library downstream of
Mathlib \citep{tauceti2026}.

\section{Formalization process}
\label{sec:system}

The workflow in \cref{fig:workflow} can be summarized by five recurring
activities.  ``Gather Context'' corresponds to the inputs at the left of the
figure, while ``Formalize'' below includes its compile-and-revise cycle:

\begin{itemize}
\setlength{\itemsep}{0.25em}
\item \textbf{Gather Context.}  We collected the relevant PDFs, built
  transcriptions of the source passages, and wrote prose distillations of the
  definitions, assumptions, and arguments needed for the target result.  These
  materials fixed the mathematical context against which later Lean statements
  were reviewed.

\item \textbf{Decompose.}  We worked out a sequence of intermediate results and
  dependencies that would support the target theorem.  This serves a role similar
  to the blueprints used in other large Lean formalizations
  \citep{liquidblueprint2023,pfrblueprint2024,zhu2026leanarchitect}: it gives the
  implementation agents a mathematical route through a task larger than one
  proof request.

\item \textbf{Find Foundations.}  Before rebuilding an intermediate result, we
  searched Mathlib, existing Lean developments, and the literature.  In this
  project, for example, Spectra supplied spectral and operator-theoretic material
  that we reused or extended \citep{bornemann2026spectra}.  Missing foundations
  were developed locally for reuse by later targets.

\item \textbf{Formalize.}  Agents implemented the resulting tasks in Lean,
  compiled them, and revised definitions and proofs in response to compiler
  feedback.  We used ChatGPT and Claude Code across this work rather than tying
  a fixed role to a particular system.

\item \textbf{Skeptical Review.}  When possible, an agent that had not implemented
  the result compared the source passage with the resulting theorem statement.
  The review checked hypotheses, objects, scalar field, dimension, bounded or
  unbounded operator scope, norm class, gap assumptions, constants, and
  directionality.  Implementing agents sometimes declared a target complete
  before this comparison; reviewers could also demand a stronger statement than
  the source required.  The human authors resolved such disagreements by
  returning to the source mathematics.
\end{itemize}

The stages loop.  A mismatch can require more source context, a different
decomposition, additional foundations, or a revised formal statement.  As the
formalization grew, we periodically refactored definitions and intermediate
results so foundations developed for one target could be reused by later ones.
We also added project-specific tools for querying compiler-derived declarations
and maintaining the accumulated source-to-statement records.

A target left the loop when Lean checked its formal treatment and a separate
source-to-statement review found its hypotheses, mathematical objects,
conclusion, and scope consistent with the source.  When that correspondence
used a non-obvious representation, the required identification also had to be
established.  Proposition~4.4 instead exits through a checked counterexample to
the printed claim together with a proved repair.  A later review can still find
a mismatch missed by an earlier pass, as in the example below.
Appendix~\ref{app:system} records the concrete project artifacts and checks.

% Historical provenance for editors (not rendered): the maintained result
% record reached 29/29 on 12 Aug. 2026, later lost that state, and again reported
% 29/29 at commit 8b3f0f39 on 17 Aug.  The ambient sin 2Theta mismatch below was
% already present in that 17 Aug. snapshot, before the Palomar-oriented cleanup
% beginning 28 Aug.  Commit 8499471d on 31 Aug. exposed it by requiring each
% printed clause to be supported by one theorem carrying the relevant result-wide
% scope.  Real and complex directed residual witnesses were already unbounded at
% the 17 Aug. checkpoint; real and complex ambient witnesses were bounded.  A
% later review also corrected which operator's blocks carry the printed spectral
% gap.  Keep this chronology and repair history in comments/appendix evidence;
% the rendered example should state the mathematical mismatch directly.
\subsection{Examples of semantically misaligned statements}
\label{sec:scope-mismatch}

Formalization~\ref{lst:sin2theta-mismatch} was accepted as evidence for the
ambient clause of the Section~2 $\sin 2\Theta$ result.  Davis--Kahan allow
unbounded self-adjoint operators, while this declaration is stated in a section
where \texttt{A B : E $\to$L[$\mathbb C$] E} are bounded continuous linear maps.
Real and complex ambient analogues were both present.  The mismatch was bounded
versus unbounded operator scope.

% Presentation note (not rendered): scripts/build_evidence.py extracts the
% historical declaration and its section-level operator context from Git and
% writes the exact text verbatim to generated/historical_scope_mismatch_exact.lean.
% The file consumed by listings is a separate presentation artifact: it carries
% a header pointing back to that exact sidecar, applies the audited
% PaperUnitaryInvariantNorm -> SymmetricNormingFunction readability rename, and
% replaces Lean Unicode notation with ASCII LeanLit... sentinels. The leanpaper
% literate= table above renders those sentinels as LaTeX glyphs. Do not audit
% source fidelity from the presentation file; use the exact sidecar/JSON
% provenance instead.
\begin{nolinenumbers}
\begin{minipage}{\linewidth}
\lstinputlisting[
  style=leanpaper,
  firstline=6,
  caption={Historical ambient $\sin 2\Theta$ theorem with bounded operator scope.},
  label={lst:sin2theta-mismatch}
]{generated/historical_scope_mismatch_presentation.lean}
\end{minipage}
\end{nolinenumbers}

The review aggregated coverage across declarations.  An unbounded directed theorem
supplied the scope, while the bounded theorem above supplied the ambient inequality.
No single declaration established the ambient claim at the source's unbounded
scope.  The revised check pairs each source clause with a theorem carrying its
applicable scope.

Two further repairs exposed different statement-level discrepancies:
\begin{itemize}
\setlength{\itemsep}{0.15em}
\item In the Section~2 $\sin\Theta$ theorem, the earlier source-facing statement
used a stronger norm abstraction, required residual ideal membership, and concluded
that the angle operator belongs to the norm's ideal.  The repaired boundary uses the source's
where-defined normalized unitarily invariant norm semantics.
\item Theorem~8.1(ii),(iii) was stated on blocks extended by zero to the ambient
space, with ambient-size indexing and gauges.  The repaired finite-dimensional
statements use the ordered eigenvalues of the blocks themselves and the block
dimension.
\end{itemize}

\paragraph{Current source-facing sine-theta statement.}
For $\mathbb K\in\{\mathbb R,\mathbb C\}$ on a separable Hilbert space, under
the source's self-adjointness, residual, exact-decomposition, and spectral-gap
hypotheses (including its unbounded-operator scope), the current endpoints state
for every normalized unitarily invariant norm $N$:
\begin{equation}
  N(\sin\Theta_0)\ \text{defined}\;\land\;N(R)\ \text{defined}
  \quad\Longrightarrow\quad
  \delta N(\sin\Theta_0)\le N(R),
  \label{eq:source-facing-sin-theta}
\end{equation}
where $P$ is the orthogonal projector onto the exact subspace,
$\sin\Theta_0=(I-P)E_0$, and $R=AE_0-E_0A_0$.  The proof uses a stronger
symmetric-norming theorem, retained separately as a generalization.

% Candidate future example of review overreach (not rendered): Theorem 8.1(ii)
% ends with "and natural infinite-dimensional extensions."  A prior review
% selected approximation-number extensions as the source meaning.  The paper
% does not choose a unique extension, so those theorems are retained as
% generalizations rather than source-equivalent evidence.  If this is used in
% the manuscript later, present it as a calibrated-review example, not as repair
% chronology.

\section{Conclusion}

\paragraph{Did we formalize Davis--Kahan?}
All 29 tracked results have checked treatments: 28 Lean proofs and a
finite-dimensional counterexample plus proved repair for Proposition~4.4
(Appendix~\ref{app:data}).  The current source-facing statements incorporate the
repairs described above, but we have no calibrated stopping rule or certificate
of complete source correspondence.

Lean establishes the encoded proposition.  Source review asks whether that
proposition preserves the intended hypotheses, objects, and scope without
exceeding what the source states; how to validate and stop such review remains
open.

\section*{Tool and computational resource disclosure}
Following the Leiden Declaration's disclosure recommendation
\citep{leiden2026declaration}, ChatGPT and Claude Code were substantially
used in the writing process, including drafting and revision, literature
discovery, and organization of the practitioner survey.  The same systems were
used for formalization planning, Lean implementation, and source-to-signature
review as described in \cref{sec:system}.  Lean~4 and Mathlib provided formal
checking.  Human authors selected the research questions and claims,
adjudicated source interpretations, reviewed and revised the manuscript, and
retain responsibility for its contents and citations.

\begin{ack}
This material is based upon work supported by the Defense Advanced Research
Projects Agency (DARPA) under Contract No. HR001125CE017. Any opinions, findings,
and conclusions or recommendations expressed in this material are those of the
author(s) and do not necessarily reflect the views of the Defense Advanced
Research Projects Agency (DARPA).
\end{ack}

\clearpage
\begingroup
\sloppy
\bibliographystyle{unsrtnat}
\bibliography{references}
\endgroup

\clearpage
\appendix
\section{Supplementary data}
\label{app:data}

The supplementary tables contain \PractitionerAccountCount{} public
first-person practitioner accounts supported by \PractitionerSourceCount{}
practitioner-source records.  Known repeated practitioners form
\PractitionerOverlapClusterCount{} overlap clusters; one controlled human--AI
workflow study is kept separate.  Accounts are project/workflow episodes rather
than unique-person records, and multiple sources may support one account.  The
collection was assembled through LLM-assisted web search and is not a systematic
or representative sample.

The screening table records \PractitionerScreenedCandidateCount{} candidates and
their dispositions.  It begins on 8 September 2026 and does not reconstruct
every source encountered in earlier searches.

The supplementary material also contains the 29-result register used for the
Davis--Kahan formalization.  It counts the four unnumbered headline theorems in
Davis--Kahan Section~2 and every named theorem, proposition, lemma, and corollary
established in the source.  Definitions, proof steps, examples, immediate
consequences, and open or deferred claims are not counted separately.
Davis--Kahan Proposition~4.4 remains a counted result because its formal
treatment gives a counterexample to the printed claim and proves a repaired
theorem.

\ifshowartifacturl
Project materials: \href{\artifacturl}{public GitHub repository}.
\else
Identifying repository URLs are suppressed in this double-blind build and are
included in the supplementary artifact.
\fi

\section{Formalization procedure}
\label{app:system}

\Cref{sec:system} summarizes the workflow.  The project used a pinned
Lean/Mathlib environment together with reusable local foundations.  Source
transcriptions, prose summaries, result records, and review judgments were kept
outside individual model conversations so work could continue across sessions.
The reviewed snapshot specifies Lean~4.34.0-rc1, Mathlib revision
\texttt{72f9c607bb3a}, and Tau Ceti revision \texttt{1b39d420ac84}.
The reusable mathematics package is \texttt{ForTauCeti}; source-facing
Davis--Kahan declarations live in the separate \texttt{DavisKahan} library.
Full dependency revisions are recorded in the supplied Lake manifest.
These identify the reproduction environment; the proposed replacement below
has not yet been validated in it.
As the project grew, we also developed tools for inspecting elaborated
declarations and checking that recorded formal evidence still referred to the
current development. Reusable foundations lived below the paper-facing theorem
modules. Agent tasks named an intended statement and its source context; agents
returned Lean declarations and obstruction notes. Compilation, axiom inspection,
and signature extraction supplied mechanical evidence, while correspondence
judgments were recorded separately. Source review included inherited section
variables and the fields of structures occurring in a theorem type, because
both can restrict a claim without appearing in its concluding inequality.

Source review compared each tracked Davis--Kahan result with the elaborated Lean
statement and definitions that affected its meaning.  Recurring questions were
the scalar field, dimension, bounded or unbounded operator scope, norm class,
gap conditions, constants, directionality, and the identity of angle and
residual operators.  Mechanical checks established properties of the formal
artifact; they did not determine whether our reading of the source was correct.

ChatGPT and Claude were used for implementation, mathematical discussion,
search, and review.  Human supervision consisted mainly of monitoring progress,
supplying framing, asking models to explain inconsistencies, and requesting
additional or independent reviews.  The humans directing the formalization were
not Davis--Kahan subject-matter experts, so disagreements about source meaning
were often investigated by probing the models in different ways rather than by
an independent expert reading of the theorem.

\section{Davis--Kahan background and Formalization 1}
\label{app:formalization1}

In the source notation, $T=A+H$ is the perturbed self-adjoint operator,
$P\mathcal H$ reduces $A$, and $Q\mathcal H$ reduces $T$
\citep{davis1970rotation}. With $E_0$ an isometry onto $P\mathcal H$,
$AE_0=E_0A_0$ defines the unperturbed trial restriction. The general residual is
\begin{equation}
 R=TE_0-E_0A_0=HE_0.
 \label{eq:dk-general-residual}
\end{equation}
In finite dimensions, the Rayleigh--Ritz choice is
\begin{equation}
 A_0=E_0^*TE_0,
 \label{eq:dk-trial-compression}
\end{equation}
which makes the trial-side diagonal block $H_0$ of the perturbation zero.
Only under this specialization is the residual exactly the component outside
the trial space:
\begin{equation}
 R=(I-P)TE_0.
 \label{eq:dk-trial-residual}
\end{equation}
It then vanishes precisely when the trial space is invariant under $T$.
In general, $R=0$ means that the specified trial operator and embedding satisfy
the exact intertwining relation, which is stronger than invariance alone.

Formalization~1 displays an \emph{ambient} theorem, not a residual theorem.
Its local names \texttt{A} and \texttt{B} denote the two bounded self-adjoint
operators; \texttt{U} reduces \texttt{A} and \texttt{V} reduces \texttt{B}.
The perturbation in its conclusion is \texttt{B - A}. No trial operator
\texttt{M} or residual \texttt{R} occurs in that declaration.

The hypotheses \texttt{hUspec} and \texttt{hUspec'} separate the restrictions
of \texttt{A} to $U$ and $U^\perp$ by a gap $d>0$: their spectra lie in
$[a,b]$ and outside $(a-d,b+d)$, respectively. This gap is on the first
operator. For the ambient double-angle bound, Davis--Kahan explicitly permit
interchanging the two operators. Thus this placement is not a defect in the
displayed ambient theorem. The directed residual inequality does not admit
that interchange in general. The historical discrepancy is the bounded
ambient operator type, not the choice of which operator supplies this gap.

For two subspaces of the same finite dimension with orthonormal coordinate
maps $E_0$ and $F_0$, their principal angles $\theta_j\in[0,\pi/2]$ satisfy
$\cos\theta_j=\sigma_j(F_0^*E_0)$.  Equivalently, they pair directions in the two
subspaces by their overlap.  The directed sine measures the part of a trial
direction outside the exact subspace; its matrix representative is
$(I-Q)E_0$.  The doubled-angle quantity applies $\sin(2\theta_j)$ to the
angles, with the trial-to-exact orientation fixed by the source.  The bound
\[
  \delta\,N(\sin 2\Theta_0)\le 2N(R)
\]
therefore controls the doubled-angle quantity by the residual relative to the
gap.  A small doubled sine alone does not imply a small angle:
$\sin(2\theta)$ vanishes at both $0$ and $\pi/2$.  Recovering a small-angle bound
requires an appropriate angle branch, such as $0\le\theta\le\pi/4$.

Formalization~1 states the estimate for complex Hilbert spaces.  Its ambient
Hilbert-space assumptions are omitted from the displayed code.  In the syntax
shown there, \texttt{E }$\mathrel{\to}\!\mathtt{L}[\mathbb C]$\texttt{ E}
means a bounded continuous complex-linear map $E\to E$.  It elaborates to
\texttt{ContinuousLinearMap (RingHom.id Complex) E E}; its boundedness is the historical scope restriction. The complex field
is not itself the project-level mismatch: a real sibling was present.  The following
table uses the same surface syntax as the displayed formalization.

\begin{table}[H]
\centering
\small
\caption{Reader's guide to the displayed ambient theorem in Formalization~1.}
\label{tab:formalization1-glossary}
\begin{tabularx}{\linewidth}{@{}>{\raggedright\arraybackslash}p{0.39\linewidth}>{\raggedright\arraybackslash}X@{}}
\toprule
Lean expression & Mathematical role \\
\midrule
\texttt{A B : E }$\mathrel{\to}\!\mathtt{L}[\mathbb C]$\texttt{ E}
  & Both ambient operators are bounded continuous complex-linear maps. \\
\texttt{IsSelfAdjoint A}
  & $A=A^*$; likewise for $B$. \\
\texttt{Reduces A U}
  & Both $U$ and $U^\perp$ are invariant under $A$; the other hypothesis
    says that $V$ reduces $B$. \\
\texttt{U.HasOrthogonalProjection}
  & The orthogonal projection onto $U$ is available; likewise for $V$.
    These instances are supplied by the surrounding context. \\
\texttt{compressOperator U A}
  & The compression of $A$ to $U$, equal to its restriction because $U$
    reduces $A$. \\
\texttt{spectrum }$\mathbb R$\texttt{ (...)}
  & The real spectrum of that self-adjoint compression. \\
\texttt{SymmetricNormingFunction}
  & The presentation's audited rename of the historical norm record.
    It supplies the ideal and gauge; the exact sidecar retains the original name. \\
\texttt{N.Mem (B - A)}
  & The perturbation lies in the ideal on which $N$ is finite. Unlike the
    current where-defined boundary, this historical theorem then \emph{concludes}
    angle ideal membership as well. \\
\texttt{N.gauge (B - A)}
  & The real-valued perturbation norm used on the right-hand side. \\
\texttt{paperSinTwoAngleOperatorC U V}
  & The ambient doubled-angle operator. Its two angle blocks must not be
    replaced by a single directed block when comparing general ideal norms. \\
\bottomrule
\end{tabularx}
\end{table}

The historical ambient conclusion asserts ideal membership and bounds the
ambient doubled-angle gauge by twice the gauge of \texttt{B - A}.  These norm and spectral conditions
are part of the statement.  The displayed declaration is complex and bounded;
a real ambient sibling also existed.  The project-level mismatch was therefore
the bounded ambient operator type, which did not cover Davis--Kahan's unbounded
setting.

\section{Common-domain double-angle replacement}
\label{app:domain-repair}

Let $A$ and $T$ be self-adjoint operators on the same dense domain
$\mathcal D\subseteq\mathcal H$. Let $P$ reduce $A$, let $Q$ reduce $T$,
and place the source gap on $T|_{Q\mathcal H}$ and $T|_{Q^\perp\mathcal H}$.
The directed clause takes a bounded map $R:P\mathcal H\to\mathcal H$ such that
\[
 Tp=Ap+Rp \qquad(p\in P\mathcal H\cap\mathcal D).
\]
The trial restriction is $A_0=A|_{P\mathcal H}$, with its own dense domain.
Since $A$ and $T$ have the same domain, $T\iota_P$ and $\iota_PA_0$ have
exactly that common domain. Neither the trial restriction nor $T-A$ is required
to be bounded. This gives the source's directed residual setup.
The ambient clause separately takes a bounded self-adjoint $H$ with $T=A+H$.
A residual hypothesis is unnecessary for that clause.

The new proof code reduces the domain issue to a bounded off-diagonal block.
Write $P^\perp=I-P$, let $\iota_P:P\mathcal H\hookrightarrow\mathcal H$, and set
\[
 C=P^\perp R\iota_P^*,\qquad S=C+C^*,\qquad J=2P-I.
\]
Reduction of $A$ and the common-domain equality imply that $P$ and $J$ preserve
$\mathcal D$. For $x\in\mathcal D$ the residual identity gives
$Cx=P^\perp TPx$. Self-adjointness of $T$, tested on the dense domain, gives
$C^*x=PTP^\perp x$. Thus
\[
 (C-C^*)x=TPx-PTx,\qquad (T-2S)Jx=JTx.
\]
The existing reflection estimate then bounds the doubled-angle block in every
Ky Fan norm by $2\|C\|_k\le 2\|R\|_k$. The existing angle correspondence and
the norm record's where-defined Fan law give the directed inequality in
\cref{eq:sin2theta}. The ambient clause uses the separate existing ambient
theorem. The proposed Lean module follows these steps; its compiler validation
remains pending. These equations explain the proposed repair, not a claim that
the new Lean declarations have already been checked.

The restriction in the current theorem is also visible with a finite gap
interval. Take $A=T=D\oplus0$ on $\ell^2\oplus\ell^2$, with
$D(x_n)=((n+1)x_n)$, let $P$ select the first summand and $Q$ the second.
Then $\sigma(T|_Q)=\{0\}$, $\sigma(T|_{Q^\perp})\subseteq[1,\infty)$,
and the gap condition holds with $[\beta,\alpha]=[0,0]$, $\delta=1$.
The source's dimension conditions hold and $H=R=0$. The angle is $\pi/2$,
so its doubled sine is zero; $P\mathcal H$ still fails the whole-space domain
hypothesis. This example also shows why a small doubled sine alone does not
select the small-angle branch.

The residual-only boundedness issue is independent. With
$A=(-D)\oplus D$, $T=(-D)\oplus2D$, and $P=Q$ the first summand,
the operators have the same domain, $R=0$, and the perturbed spectral parts
are separated. The difference $T-A=0\oplus D$ is unbounded. The directed
source inequality remains applicable, although an endpoint requiring a
bounded ambient perturbation cannot be instantiated.

\paragraph{Norm representation.}
The current norm record contains a symmetric ideal gauge, rank-one
normalization, and the comparison implication
\[
 \bigl(\forall k,\ \|X\|_k\le\|Y\|_k\bigr)
 \ \Longrightarrow\ N(X)\le N(Y)
\]
when both $X$ and $Y$ belong to the norm's domain. The theorem exposing this
property projects the corresponding structure field. This uses the
source-cited Fan comparison as packaged mathematical background; it is not an
independent formal derivation of Fan dominance from the bare norm axioms.
The new common-domain code retains that boundary and makes no additional
norm-representation claim.

\paragraph{Transcription check.}
The supplied modernized transcription states the factor $2$ in the headline
residual inequality but omits it in the residual step of the Section~7 proof.
The factor $2$ is necessary. For
\[
 A=0,\quad T=H=\begin{pmatrix}0&1\\1&0\end{pmatrix},\quad
 P\mathbb R^2=\operatorname{span}(e_1),\quad
 Q\mathbb R^2=\operatorname{span}(e_1+e_2),
\]
the spectral gap is $2$, the angle is $\pi/4$, and $\|R\|=1$.
Hence $\delta\|\sin2\Theta_0\|=2=2\|R\|$; the inequality with factor $1$
is false. The reflection calculation also gives $H-JHJ$ equal to twice the
off-diagonal part. This identifies an inconsistency in the supplied
transcription; locating its origin requires comparison with the printed scan.

\section{Lean publication activity}

\begin{figure}[H]
  \centering
  \input{generated/lean_activity_timeline_tikz.tex}
  \caption{Monthly publication counts reproduced from the Papers With Lean
  public statistics series, January 2025 through August 2026
  \citep{paperswithlean2026stats}.  Counts are grouped by the captured data's
  \texttt{published} month, and only complete months are shown.  The fraction
  that used AI is unknown.}
  \label{fig:lean-activity}
\end{figure}



\section{Semantic-review page}
\label{app:review-page}

\IfFileExists{figures/semantic-alignment-sine-theta-row.png}{
\begin{figure}[H]
  \centering
  \includegraphics[width=\linewidth,trim=0 250 0 0,clip]{figures/semantic-alignment-sine-theta-row.png}
  \caption{A review interface developed late in the project.  The source passage
  and inherited context appear at left, the claimed source-to-Lean
  correspondence in the middle, and compiler-derived Lean evidence at right.
  Most reviews reported in the main paper used LLM source/signature comparison;
  the interface was added later.}
  \label{fig:dashboard}
\end{figure}

}{
The dashboard image is omitted from this review build pending regeneration.
It is not used as evidence for the current source-correspondence judgment.
}

The interface consolidates source passages, correspondence judgments, and
compiler-derived signatures previously inspected through files and Lean
queries.  Understanding a signature can also require the definitions of its
structures, predicates, coercions, and typeclass assumptions.  Lean Compass's
semantic dependency view suggests how these definitions could be exposed while
omitting proof-only dependencies \citep{yanahama2026leanatlas}.

\section{Resource utilization}
\label{app:resources}

To estimate the cost of the formalization effort, we maintained a resource
accounting record.  The compact snapshot below has an accounting cutoff of
13 August 2026.  Retained model turns and token quantities are incomplete
measured lower bounds; price and serving-energy rows are estimates derived from
those observations.  The energy estimate excludes training, embodied hardware,
client devices, and network energy.

\begin{table}[H]
\centering
\scriptsize
\caption{Observed lower bounds and derived estimates at the 13 August 2026
accounting cutoff.  These are not project totals.}
\begin{tabularx}{\linewidth}{@{}p{0.28\linewidth}p{0.17\linewidth}>{\raggedright\arraybackslash}X@{}}
\toprule
Quantity & Lower bound / estimate & Basis / interpretation \\
\midrule
\input{generated/resource_snapshot_table.tex}
\end{tabularx}
\end{table}


\begin{table}[H]
\centering
\scriptsize
\caption{AI systems and model families used during the formalization.
Token telemetry is incomplete.}
\label{tab:formalization-systems}
\begin{tabularx}{\linewidth}{@{}p{0.24\linewidth}>{\raggedright\arraybackslash}X@{}}
\toprule
System & Model families observed \\
\midrule
\input{generated/model_systems_table.tex}
\end{tabularx}
\end{table}

\begin{table}[H]
\centering
\scriptsize
\caption{Recovered model-resolved token measurements at the 13 August 2026
accounting cutoff.  These are measured lower bounds for the five labels with
recovered telemetry; models in \cref{tab:formalization-systems} that do not
appear here were still used.}
\label{tab:model-tokens-appendix}
\resizebox{\linewidth}{!}{% Generated by scripts/build_accounting.py; do not edit by hand.
\begin{tabular}{lrrrr}
\toprule
Model & \multicolumn{4}{c}{Retained measured tokens} \\
 & Input & Cache write & Cache read & Output \\
\midrule
claude-opus-5 & 140,093 & 116,791,540 & 17,395,271,984 & 34,338,937 \\
claude-opus-4-8 & 693,366 & 58,990,139 & 3,744,966,700 & 15,967,324 \\
claude-fable-5 & 380,759 & 27,789,153 & 1,512,917,593 & 7,985,843 \\
gpt-5.6-sol & 6,502,664 & 0 & 343,618,944 & 635,601 \\
gpt-5.6-terra & 3,798,915 & 0 & 382,951,936 & 563,722 \\
claude-sonnet-5 & 78 & 110,211 & 2,256,550 & 26,074 \\
\bottomrule
\end{tabular}
}
\end{table}

Historical model labels are retained verbatim, with exact recovered telemetry used
when it is available.


\clearpage
\input{checklist.tex}

\end{document}
